\documentclass[10pt,a4paper]{amsart}
\usepackage[margin=19mm]{geometry}
\usepackage{amsmath,amssymb,mathtools}
\usepackage[scr=rsfs,cal=euler]{mathalfa}
\usepackage{xfrac}
\usepackage[unicode,hidelinks,bookmarks=false]{hyperref}
\newcommand{\ie}{i.e.\ }
\newcommand{\cf}{cf.\ }
\newcommand{\resp}{respectively\ }
\newcommand{\Subsection}{\subsection}
\providecommand{\nobreakdash}{\nobreak\hbox{-}}
\setlength{\emergencystretch}{2em}
\multlinegap0pt
\usepackage{bm}
\usepackage{bbm}
\usepackage{dsfont}
\usepackage{xspace}
\usepackage{cancel}
\usepackage{mathtools}
\usepackage{amsthm}
\makeatletter
\@ifundefined{theorem}{\newtheorem{theorem}{Theorem}}{}
\@ifundefined{lemma}{\newtheorem{lemma}[theorem]{Lemma}}{}
\@ifundefined{proposition}{\newtheorem{proposition}[theorem]{Proposition}}{}
\makeatother
\makeatletter
\expandafter\ifx\csname AMS\endcsname\relax
\newenvironment{AMS}{\noindent \textbf{AMS subject classification.}}{}
\fi
\newcommand{\KL}[2]{\text{KL}(#1\,|\,#2)}
\def\P{\mathcal P}
\def\R{\mathbb R}
\DeclareMathOperator*{\amin}{argmin}
\renewcommand{\d}{\mathrm{d}}
\newcommand{\ddt}{\frac{\d}{\d t}}
\expandafter\ifx\csname keywords\endcsname\relax
\newenvironment{keywords}{\noindent \textbf{Keywords.}}{}
\fi
\makeatother
\title[Computing Optimal Transport Plans via Min-Max Gradient Flows]{Computing Optimal Transport Plans via Min-Max Gradient Flows}
\author{Lauren Conger, Franca Hoffmann, Ricardo Baptista, Eric Mazumdar}
\date{Source: arXiv:2504.16890v2 (CC BY 4.0)}
\begin{document}
\maketitle

\noindent\emph{Source and licence.} Computing Optimal Transport Plans via Min-Max Gradient Flows. Version arXiv:2504.16890v2 (CC BY 4.0), distributed under the Creative Commons Attribution 4.0 International licence, as recorded in the arXiv licence metadata for this exact version and shown on the abstract page https://arxiv.org/abs/2504.16890v2. Full rights notice: licenses/arxiv-2504.16890-CC-BY-4.0.txt.\par\smallskip

\noindent\textit{Editorial scope.} Counted unit: the Proposition whose source label is \texttt{prop:Lambda-properties} in the article (source0.tex lines 885--887, source label \texttt{prop:Lambda-properties}), delivered together with the article's own complete original proof of it (source0.tex lines 888--912, one \texttt{proof} environment of 25 lines). No mathematics is written, repaired or re-derived: statement and proof are the article's own text line for line. Required context retained uncounted: eq:OT\_problem (lines 538--540); eq:best\_response\_dynamics (lines 579--584); def:V (lines 676--678); lem:best-response-marginals (lines 708--719); lem:continuity\_V (lines 797--799); thm:Lambda-upper (lines 838--844); lem:V-Lip (lines 871--876). Every definition, display and label of those lines is retained unchanged. Omitted from this package and not redistributed: 1--537, 541--578, 585--675, 679--707, 720--796, 800--837, 845--870, 877--884, 913--1062. Nothing inside a retained excerpt is omitted or altered: every retained body is delivered exactly as the article's lines are written, so no omission receipt of any kind accompanies this package. Citations: the retained excerpts carry no citation command at all, so no bibliography entry of the article is transcribed and no reference is redistributed. Reference adaptation: none was needed and none was made; every reference inside the delivered unit resolves inside it, so no unresolved reference or citation is produced. Printed slips kept literal: no printed word, symbol, hypothesis or number of the counted statement or of its proof was altered. Layout and class: the article's own class is \texttt{article} with packages including authblk, silence, fontenc, inputenc, csquotes, babel, paralist, enumitem, geometry, xcolor, graphicx, bbm, amsmath, amssymb; the wrapper is the lane's pinned \texttt{amsart} editorial wrapper at 10pt on A4, which restates the theorem-like environments the retained text needs and adds the article's own macro definitions that the retained text needs (\texttt{\textbackslash KL}, \texttt{\textbackslash P}, \texttt{\textbackslash R}, \texttt{\textbackslash amin}, \texttt{\textbackslash d}, \texttt{\textbackslash ddt}), with extra wrapper packages (bm, bbm, dsfont, xspace, cancel, mathtools). The delivered numbering is the unit's own. Assets: no retained span contains a figure environment, a graphics-inclusion command, a PDF-page inclusion, a table or a TikZ picture; no image file is copied into this package, the package declares no asset dependency and no image file is redistributed with it. Licence: version arXiv:2504.16890v2 of this article is distributed under the Creative Commons Attribution 4.0 International licence, as recorded in the arXiv licence metadata for this exact version and shown on the abstract page https://arxiv.org/abs/2504.16890v2. A full rights notice is delivered at licenses/arxiv-2504.16890-CC-BY-4.0.txt. Characters: every retained excerpt is pure ASCII, so the delivered engine, XeLaTeX with its default Latin Modern text font, reports no missing character.
\begin{align}\label{eq:OT_problem}
    \inf_{\rho\in\Gamma(\mu,\nu)} \iint c(x,y)\d \rho(x,y)\,.
\end{align}

\begin{align}\label{eq:best_response_dynamics}
\begin{split}
        \ddt \Lambda(t) &= \partial_\Lambda  E(b[\Lambda],\Lambda;\mu,\nu) \\
    b[\Lambda] &\coloneqq \amin_{\rho\in\P_2} E(\rho,\Lambda;\mu,\nu) \,.
\end{split}
\end{align}

\begin{align}\label{def:V}
   V(\Lambda):=\KL{b_1[\Lambda]}{\mu} + \KL{b_2[\Lambda]}{\nu}\,,
\end{align}

\begin{lemma}[Marginals of best response]\label{lem:best-response-marginals}
Fix $\Lambda>0$ and denote $Z[\Lambda]:= \iint e^{-c(x,y)/\Lambda}\d\mu(x)\d\nu(y)$. Then the product of the marginals of the best response $b[\Lambda]$ satisfies
\begin{align}\label{eq:EL}
    b_1[\Lambda](x)b_2[\Lambda](y)
    =\frac{1}{Z[\Lambda]}e^{-c(x,y)/\Lambda}\mu(x)\nu(y)=:\sigma[\Lambda](x,y)\,,
\end{align}
where the normalization constant $Z=Z[\Lambda]$ is bounded,
\begin{align}\label{eq:Z}
    0<e^{-c^*/\Lambda} \le Z[\Lambda]\le 1 \quad \forall \,\Lambda>0\,,
\end{align}
with $c^*$ the optimal value of \eqref{eq:OT_problem}.
\end{lemma}

\begin{lemma}[Continuity of $V$]\label{lem:continuity_V}
    $V:(0,\infty)\to[0,\infty)$ defined in \eqref{def:V} is continuous.
\end{lemma}

\begin{theorem}[Regularization growth condition]\label{thm:Lambda-upper}
    Under the dynamics \eqref{eq:best_response_dynamics}, the regularization parameter $\Lambda$ satisfies
    \begin{align*}
        &\Lambda(t) \le \sqrt{2(c^* t + c_0)} \quad \forall t\ge 0\,,
    \end{align*}
    where $c^*<\infty$ is the optimal value of \eqref{eq:OT_problem} and $c_0>0$ is a constant dependent on $\Lambda(0)>0$.
\end{theorem}

\begin{lemma}[Lipschitzness of $V$]\label{lem:V-Lip}
There exists a constant $0<c_0<\infty$ such that 
$$
\left\|\frac{\d V}{\d\Lambda}\right\|_\infty=\sup_{\Lambda\in[0,\infty)}\left|\frac{\d V}{\d\Lambda}(\Lambda)\right|<c_0\,.
$$
\end{lemma}

\begin{proposition}[Properties $\Lambda(t)$]\label{prop:Lambda-properties}
    The ODE $\dot\Lambda=V(\Lambda)$ with $\Lambda(0)>0$ has a unique solution $\Lambda(\cdot)\in C([0,\infty))$ that satisfies $V(\Lambda(t))\to 0$ as $t\to\infty$. In other words, the best response $b[\Lambda(t)]$ satisfies the marginal constraints asymptotically.
\end{proposition}

\begin{proof}
  For $\Lambda(0)>0$, define $C_0:= \frac{c^*}{\Lambda_0(\Lambda_0+1)}$. Since $\dot\Lambda=V(\Lambda)\ge 0$, we have 
$$
\Lambda(t)\ge \Lambda(0)= -\frac{1}{2}+\sqrt{\frac{1}{4}+\frac{c^*}{C_0}}\,\Rightarrow\,
V(\Lambda)\le \frac{c^*}{\Lambda}\le C_0(1+\Lambda)\,.
$$
for all $t\ge 0$. Thanks to continuity of $V$ (Lemma~\ref{lem:continuity_V}), Lipschitzness of $V$ (Lemma~\ref{lem:V-Lip}) and the growth condition from Theorem~\ref{thm:Lambda-upper}, we obtain global existence and uniqueness of a continuous solution $\Lambda(t)$ to the ODE $\dot\Lambda=V(\Lambda)$.
% with $0<\Lambda(0)\le\Lambda(t)\le \sqrt{2(c^* t + c_0)}$.  
Further, $\Lambda(t)$ is non-decreasing, so it must either remain bounded for all times, or diverge to $+\infty$. Thus, we can write $\Lambda(t)\to\Lambda_\infty\in [\Lambda(0),\infty]$, and by continuity, $$V(\Lambda(t))\to V_\infty:=V(\Lambda_\infty)\in [0,\infty]\,.$$ 
Note from Lemma~\ref{lem:best-response-marginals}  for all $t\ge 0$,
$$
V(\Lambda(t))= \frac{1}{\Lambda(t)}\left(E_d(\Lambda(t))-\int c\d b[\Lambda(t)]\right)\le \frac{c^*}{\Lambda(0)}\,,
$$
and so $V$ is uniformly bounded from above. Thus, 
there are two possibilities: either (a) $V_\infty\in (0,\infty)$, or (b) $V_\infty=0$.
In Case (a), there exists $T>0$ such that for all $t\ge T$,
\begin{align*}
    \dot\Lambda(t)=V(\Lambda(t))> \frac{V_\infty}{2}>0\,.
\end{align*}
Integrating the ODE, we obtain the lower bound
\begin{align*}
    \Lambda(t)>\frac{\Lambda_\infty}{2} t + V_0 \quad \forall t\ge T\,,
\end{align*}
where $V_0=V_0(T)\in\R$. For large enough $t\ge T$, this is a contradiction to the upper bound in Theorem~\ref{thm:Lambda-upper}, so (a) cannot happen. We conclude $V(\Lambda(t))\to 0$.  
\end{proof}

\end{document}
